% Clase 20 --- Los experimentos de inducción: el anillo de Thomson (§3).
% "Anillo de aluminio sobre una bobina con CA y núcleo de hierro levita;
%  cortado, no levita; sin núcleo, tampoco."
% Tres versiones del MISMO montaje: cada panel saca un ingrediente y el efecto
% desaparece. Eso es lo que identifica qué hace falta para que haya inducción.
\begin{center}
\begin{tikzpicture}[scale=1]

  \providecommand{\bobinaCA}[1]{%
    \draw[figline, line width=1.1pt] (-0.75,-1.2) rectangle (0.75,0.55);
    \foreach \yy in {-1.05,-0.75,-0.45,-0.15,0.15,0.45} {
      \draw[figline, line width=0.8pt] (-0.75,\yy) -- (0.75,\yy);
    }
    \ifnum#1=1
      \fill[figamber!40] (-0.28,-1.2) rectangle (0.28,1.35);
      \draw[figamber, line width=0.9pt] (-0.28,-1.2) rectangle (0.28,1.35);
    \fi
    \node[figlblaux, anchor=north] at (0,-1.3) {$\sim$ CA};}

  % =============== (a) con núcleo: levita ===============
  \begin{scope}
    \bobinaCA{1}
    \draw[figline, line width=1.5pt] (-0.62,1.75) -- (0.62,1.75);
    \draw[figline, line width=1.5pt] (-0.62,1.9) -- (0.62,1.9);
    \node[figlbl, anchor=west] at (0.7,1.82) {anillo};
    \draw[figvecres, line width=1.2pt] (0,2.15) -- (0,2.85);
    \node[figlbl, figred, anchor=south] at (0,2.9) {salta};
    \node[figlbl, anchor=north, align=center] at (0,-2.1)
      {(a) anillo cerrado\\[-0.2em] con núcleo: \textbf{levita}};
  \end{scope}

  % =============== (b) anillo cortado ===============
  \begin{scope}[xshift=4.6cm]
    \bobinaCA{1}
    \draw[figline, line width=1.5pt] (-0.62,1.75) -- (-0.12,1.75);
    \draw[figline, line width=1.5pt] (0.12,1.75) -- (0.62,1.75);
    \draw[figline, line width=1.5pt] (-0.62,1.9) -- (-0.12,1.9);
    \draw[figline, line width=1.5pt] (0.12,1.9) -- (0.62,1.9);
    \node[figlbl, figred, anchor=west] at (0.7,1.82) {corte};
    \node[figlblaux, anchor=south] at (0,2.3) {no pasa nada};
    \node[figlbl, anchor=north, align=center] at (0,-2.1)
      {(b) anillo cortado:\\[-0.2em] no puede circular corriente};
  \end{scope}

  % =============== (c) sin núcleo ===============
  \begin{scope}[xshift=9.2cm]
    \bobinaCA{0}
    \draw[figline, line width=1.5pt] (-0.62,1.75) -- (0.62,1.75);
    \draw[figline, line width=1.5pt] (-0.62,1.9) -- (0.62,1.9);
    \node[figlblaux, anchor=south] at (0,2.3) {apenas se mueve};
    \node[figlbl, anchor=north, align=center] at (0,-2.1)
      {(c) sin núcleo:\\[-0.2em] el flujo es muy chico};
  \end{scope}

\end{tikzpicture}\\[0.5em]
{\small\itshape El anillo de Thomson es el experimento que mejor separa los
ingredientes, porque alcanza con sacar uno para que el efecto se caiga. Con
corriente alterna en la bobina el flujo a través del anillo varía todo el
tiempo, se induce en él una corriente y ---por Lenz--- esa corriente se opone a
la variación, de modo que el anillo es empujado hacia afuera y salta. Si se lo
corta, el circuito queda abierto: sigue habiendo FEM inducida, pero \emph{no
puede circular corriente} y no hay fuerza; ésa es la prueba de que lo que
empuja es la corriente inducida y no el campo por sí solo. Y si se quita el
núcleo de hierro el efecto se desploma, porque el hierro es lo que concentra y
multiplica el flujo que atraviesa el anillo ---el mismo $K_m$ de la primera
sección, ahora visible a simple vista---. Vale notar que acá lo que varía es el
\textbf{módulo} de $\vec B$; en el generador con imán giratorio lo que varía es
el \textbf{ángulo}. Distintas rutas, la misma ley.}
\end{center}
